

\section{Eighth Sunday after Trinity: To Worship God in Vain}


\begin{quote}

Mark 7:5–16: Then the Pharisees and the scribes asked him, “Why do your disciples not follow the tradition of the elders, but eat bread with unwashed hands?”

But he answered and said to them: "Isaiah was right when he prophesied about you hypocrites, as it is written: This people honors me with their lips, but their heart is far from me. In vain they worship me, teaching doctrines that are human commands. For you leave the commandment of God and hold to human traditions—the washing of pitchers and cups, and many other such things you do."

And he said to them: "How neatly you set aside God’s commandment so that you can keep your own traditions. For Moses said: Honor your father and your mother; and: Whoever curses father or mother must surely die. But you say: If a man says to father or mother, 'Corban,' that is, a gift to the temple, is what you might have been helped by from me, then you no longer permit him to do anything for his father or his mother, and so you nullify the word of God by the traditions you have established. And many such similar things you do."

And he called all the people to him and said to them: "Hear me, all of you, and understand. Nothing that enters a person from outside can make him unclean; but what comes out of a person is what makes him unclean. Whoever has ears to hear—let him hear."

\end{quote}

\bigskip



As may be seen from our text, the Pharisees had among other things observed that some of the Lord's disciples ate without first washing their hands. In their eyes and in the eyes of the scribes this was a great sin.

Originally it was no doubt only a custom; but it was both ancient and symbolic, and over time these two things made the custom so important and sacred that no one who would make the least claim to godliness could under any circumstances neglect it.

So much greater therefore was the astonishment and offense when a man like Jesus, who among the people had obtained such a reputation for holiness, allowed his disciples such profaneness.

The Pharisees and the scribes therefore went to him, as they had done on the occasion of the fasting, and asked him in the presence of the people: "Why do your disciples not walk according to the ordinance of the elders, but eat bread with unwashed hands?"

They likely thought that when the people heard how lightly he himself treated the sacred customs—and taught others to treat them the same way—which from childhood had been impressed upon every Israelite, they would begin to hesitate and would not cling with such admiration to this Jesus of Nazareth.

Whether this could be achieved was completely indifferent to the Lord. He did not receive honor from men. He answered them in a way that struck at their hypocrisy. The word would become for them either a savor of life unto life or of death unto death.

"The ordinances of the elders," the Lord would say, "are good enough for their use; to follow them or not follow them cannot make a person blessed. But you have misused them; you have placed them in place of God's own commandment, and therefore you worship God in vain."

He says that Isaiah's prophecy is fulfilled in them, that they are among those who honor God with their lips while their heart is far from him; that they abandon God's commandment, indeed abolish God's commandment and make it void in order to keep the ordinances of the elders and their own regulations, commandments of men and doctrines of men.

He gives an example of how they abolish, for instance, the fourth commandment—to honor father and mother—by means of their own long‑established and sacred ordinances. When someone wishes to free himself from the burden of helping or supporting an old father or mother, and he will only remain their faithful adherent and worshiper of the temple, then they know what counsel to give: "Give to the temple what you had thought to help your parents with!" It is a Corban. It is the ordinance of the elders. It gives the man a double advantage. He becomes free from the burden of his parents and receives a quiet conscience; for his gift to the temple makes him into a godly man.

But it is precisely this kind of godliness which is vain, and a severe but true word the Lord spoke of the Pharisees when he said: "You travel over land and sea to make a convert; and when he has become one, you make him a child of hell twice as much as yourselves."

And as for washing hands before eating and cleansing cups and vessels, which you are so eager about, these things can both be done and neglected without harm; but you have overturned the truth and placed also these commandments in the stead of God's commandments concerning justice and mercy and faith. For you have perverted God's truth concerning sin and have placed sin in these outward things and outward uncleanness, and you teach the people that if they only carefully wash the heathen dust from their hands and keep it carefully away from everything they enjoy, then they become true Israelites and pure before God. You forget that sin comes from within and not from without; that it is the heart and not the hands which must be washed and cleansed if a person is to become clean before God and saved. Thus you make God's word void by your ordinances and worship me in vain.

"He who has ears to hear, let him hear."

Neither the Pharisees nor the people heard. It became for them an aroma of death unto death when they rejected the Lord of glory because he told them the truth which he had heard from his Father.

But what he said to them when he walked about in the land of the Jews has since, century after century, sounded to his church on earth, and it sounds today to us.

How has this word been received by God's church?

Such doctrines as are merely human commands have in all ages been more pleasing to flesh and blood, and time after time they have caused the word of God to be abolished and made void, and men like the Pharisees of old have worshiped God in vain.

That was precisely the reason why Luther and the other reformers, reluctant as they were, had to speak against the church in which they had been baptized and brought up.

The traditions and customs of the papal church in worship and preaching had gradually become more important for holding the masses together, more useful to priestly control, and easier for people to receive than the word of God. This therefore was first set aside and then even entirely forbidden to the people.

Just as the Pharisees abolished the fourth commandment through their Corban, so the papacy had found a way to release people from all ten commandments of God, if only they paid and submitted to the commandments of the church.

And just as the Pharisees placed sin in purely outward things and forgot that from the heart proceed all manner of evil thoughts, so the papal church also placed both sin and grace entirely in outward and mechanical things—laying on of hands, consecrations, and human words—and omitted altogether the question of repentance and faith, of Christ's righteousness and the work of the Holy Spirit, until the church became half pagan and half Pharisaic, and despite prayers and kneelings and mortifications and pilgrimages and gifts of property, people nevertheless worshiped God in vain.

But these two things—diminishing the holy demands of the commandments by placing human doctrines in their stead, and placing sin in outward things, or what amounts to the same thing, making grace dependent upon other things than those which God's word has established, repentance of heart and faith—these have also faithfully followed into the Protestant church and cling like a spider’s web around every heart that seeks salvation, in order to produce a worship of God which may indeed appear splendid in its zeal, yet is nevertheless in vain.

The bitter doctrinal controversies that followed Luther in the Protestant lands—controversies that in many cases turned into religious wars in which thousands were killed—are sufficient testimony to the impure fire that burned on the church’s altar.

The spirit of unbelief and general contempt for God which filled all Protestant lands in the previous century was only the fully ripened fruit of a vain worship of God.

Also here in this land, during the many bitter doctrinal conflicts, the same spirit of carnality has revealed itself in the tendency to diminish the holy demands of repentance and the commandments and to make the way broad by such doctrines as are commandments of men, which among our people in such great measure have produced a worship of God that is vain.

When someone swears, lives in drunkenness or immorality, without daily prayer and devotion, but is zealous for the congregation and the pastor and comforts himself as to his hope of salvation on the basis of laying on of hands and the world's justification, and would readily contend and sacrifice everything for this congregational relationship, then it is a vain worship of God.

If someone is ever so zealous for his assembly and his prayer meetings, reads the Bible early and late, rebukes the world harshly for its ungodliness, and conscientiously separates himself from all who do not outwardly conduct themselves exactly as he does, but at the same time in his heart clings to the things of the world and passes by the little commandment concerning justice and love, then he worships God in vain.

When you throw yourself with all your zeal into efforts that promise an improvement in the moral or temporal condition of the peoples, but in doing so forget prayer and repentance and the daily conversion, then however much you are ready to sacrifice everything for your cause, and however much it is adorned with God's name, it is a worship of God that is vain. Of this the Lord's word is true: "These things ought you to have done, and not to leave the others undone."

And if you renounce all things, mortify your body, if you devise ever so many new forms of worship, not only in such high matters as the veneration of angels and the interpretation of prophecies, but also in small matters such as clothing, forms of speech, forms of worship, congregational regulations, and the like, but you lack love—the love which alone is poured out by the Holy Spirit in a heart that believes and daily feels its need under sin—then it is all only an empty sound, a worship of God that is vain.

To worship God rightly is to love God.

Therefore, friends: "Let us not love in word or with the tongue, but in deed and in truth."

